% ────────────────────────────────────────
%  ادامهٔ فصل ۱۲ — از جایی که قطع شد
% ────────────────────────────────────────

کردند و موج عملیات مسلحانه آغاز شد.\footnote{\en{Abrahamian, Radical Islam (1989), pp.\,218--230.}} تاریخ ۳۰ خرداد ۱۳۶۰ در تقویم جمهوری اسلامی به‌عنوان «سالروز شکست فتنهٔ منافقین» و در تقویم مجاهدین به‌عنوان «آغاز مقاومت» ثبت شده است.

\subsection{عملیات‌های مسلحانهٔ مجاهدین}

پس از ۳۰ خرداد، مجاهدین سلسله عملیات‌هایی انجام دادند که ضربات سنگینی بر حلقهٔ درونی وارد کرد:

\begin{longtable}{%
    >{\centering\arraybackslash}p{0.05\textwidth}
    >{\raggedleft\arraybackslash}p{0.14\textwidth}
    >{\raggedleft\arraybackslash}p{0.30\textwidth}
    >{\raggedleft\arraybackslash}p{0.20\textwidth}
    >{\raggedleft\arraybackslash}p{0.21\textwidth}}
\toprule
\rowcolor{cAccent!10}
\textbf{\#} & \textbf{تاریخ} & \textbf{رویداد} & \textbf{قربانیان} & \textbf{پیامد سیاسی} \\
\midrule
\endhead

۱ &
۷ تیر ۱۳۶۰ &
انفجار دفتر حزب جمهوری اسلامی &
بهشتی + ۷۲ مقام ارشد &
حذف «معمار» قانون اساسی \\
\midrule

۲ &
۸ شهریور ۱۳۶۰ &
انفجار دفتر نخست‌وزیری &
رجایی (رئیس‌جمهور) + باهنر (نخست‌وزیر) &
بی‌سابقه: رئیس‌جمهور و نخست‌وزیر هم‌زمان \\
\midrule

۳ &
پاییز ۱۳۶۰ &
ترورهای خیابانی پراکنده &
ده‌ها مقام محلی، امام جمعه‌ها &
فضای وحشت و سرکوب \\
\midrule

۴ &
۱۳۶۰–۱۳۶۲ &
عملیات‌های متعدد در شهرها &
صدها نفر (نظامی و غیرنظامی) &
اعدام‌های گسترده در پاسخ \\

\bottomrule
\caption{عملیات‌های مسلحانهٔ مجاهدین خلق (۱۳۶۰–۱۳۶۲)}
\label{tab:mek-operations}
\end{longtable}

\begin{warningBox}[انفجار ۷ تیر: لحظه‌ای تعیین‌کننده]
انفجار دفتر حزب جمهوری اسلامی در ۷ تیر ۱۳۶۰ یکی از مرگبارترین ترورهای سیاسی قرن بیستم بود: ۷۳ نفر از مقامات ارشد نظام — از جمله آیت‌الله بهشتی، رئیس دیوان عالی کشور و دبیرکل حزب — کشته شدند.\footnote{\en{Bakhash (1984), pp.\,220--225;} همچنین: \en{Moin (1999), pp.\,235--240.}}

این رویداد:
\begin{enumerate}[label=\textcolor{cAccent}{\bfseries\arabic*.}, rightmargin=2em]
    \item «معمار اصلی» ولایت فقیه در قانون اساسی را حذف کرد.
    \item موج همدردی و خشم عمومی ایجاد کرد.
    \item به بهانه‌ای برای سرکوب \keyword{بی‌سابقه} تبدیل شد.
    \item خمینی در واکنش گفت: \textit{«این‌ها [مجاهدین] بدتر از کفارند... منافق‌اند.»}\footnote{صحیفهٔ امام، ج۱۵، ص ۲۴۰. از این لحظه، واژهٔ «منافق» (= منافقین) جایگزین «مجاهدین خلق» در ادبیات رسمی شد.}
\end{enumerate}
\end{warningBox}

% ────────────────────────────────────────
\section{موج اعدام‌ها: تابستان و پاییز ۱۳۶۰}
% ────────────────────────────────────────

\subsection{آمار و ابعاد}

پس از عملیات‌های مسلحانهٔ مجاهدین، موج اعدام‌هایی بی‌سابقه آغاز شد:

\begin{itemize}[label=\textcolor{cAccent}{$\blacktriangleright$}, rightmargin=2em]
    \item بر اساس تخمین‌های مختلف، بین \keyword{۲٫۵۰۰ تا ۱۰٫۰۰۰} نفر در فاصلهٔ خرداد ۱۳۶۰ تا اسفند ۱۳۶۱ اعدام شدند.\footnote{آمار دقیق محل اختلاف شدید است. \en{Amnesty International} در گزارش ۱۹۸۲ رقم ۲٫۵۰۰ را ذکر کرد. آبراهامیان رقم بالاتری (حدود ۵٫۰۰۰) تخمین زده. مجاهدین خلق ادعای ۱۰٫۰۰۰+ دارند. \en{Abrahamian (1999), pp.\,209--228.}}
    \item اعدام‌ها فقط شامل مجاهدین نبود — اعضای چریک‌های فدایی، حزب توده، پیکار، و حتی افراد بدون وابستگی سازمانی نیز اعدام شدند.
    \item بسیاری از اعدام‌شدگان نوجوان (زیر ۱۸ سال) بودند.\footnote{\en{Amnesty International, ``Iran: Violations of Human Rights,'' 1982.}}
    \item محاکمات چند دقیقه‌ای بود: «آیا مسلمانی؟ آیا نماز می‌خوانی؟ آیا از خمینی حمایت می‌کنی؟» پاسخ منفی = اعدام.\footnote{\en{Abrahamian (1999), pp.\,210--215.} آبراهامیان بازسازی دقیقی از فرایند محاکمات ارائه داده است.}
\end{itemize}

\subsection{نقش مستقیم خمینی}

\begin{warningBox}[دستور خمینی برای اعدام‌ها]
شواهد نشان می‌دهد خمینی مستقیماً در صدور حکم اعدام‌ها دخالت داشت:

\begin{enumerate}[label=\textcolor{cAccent}{\bfseries\arabic*.}, rightmargin=2em]
    \item \textbf{فرمان عمومی:} خمینی در تیر ۱۳۶۰ فرمان «برخورد قاطع با منافقین» صادر کرد.\footnote{صحیفهٔ امام، ج۱۵، ص ۲۵۵–۲۶۰.}
    
    \item \textbf{تأیید اعدام نوجوانان:} وقتی به خمینی گزارش دادند افراد زیر ۱۸ سال اعدام می‌شوند، واکنش منفی نشان نداد.\footnote{خاطرات منتظری، فصل مربوط به ۱۳۶۰: «من به آقا [خمینی] گفتم جوان‌ها را اعدام نکنید. ایشان فرمودند: این‌ها محاربند.»}
    
    \item \textbf{انتصاب قاضیان:} دادستان‌ها و حاکمان شرع دادگاه‌های انقلاب مستقیماً توسط خمینی منصوب می‌شدند.
    
    \item \textbf{ردّ درخواست‌های رأفت:} منتظری (جانشین تعیین‌شدهٔ خمینی) بارها از خمینی خواست رأفت نشان دهد — خمینی رد کرد.\footnote{خاطرات آیت‌الله منتظری، چاپ اینترنتی، فصل‌های ۱۵–۱۸.}
\end{enumerate}
\end{warningBox}

\subsection{اعدام‌های ۱۳۶۷: «کشتار زندانیان سیاسی»}

\begin{warningBox}[فاجعهٔ تابستان ۱۳۶۷]
در تابستان ۱۳۶۷ — بلافاصله پس از پذیرش قطعنامهٔ ۵۹۸ و پایان جنگ — خمینی فرمان اعدام زندانیان سیاسی را صادر کرد. بر اساس این فرمان:

\begin{itemize}[label=\textcolor{cAccent}{$\bullet$}, rightmargin=2em]
    \item «هیئت‌های مرگ» (ترکیبی از یک قاضی شرع، یک نمایندهٔ دادستان، و یک نمایندهٔ وزارت اطلاعات) در زندان‌ها تشکیل شد.
    \item زندانیانی که سال‌ها حکم‌شان را گذرانده بودند، دوباره «محاکمه» شدند.
    \item پرسش اصلی: «آیا هنوز به سازمان [مجاهدین] وفادارید؟» پاسخ مثبت = اعدام.
    \item تخمین‌ها از ۲٫۸۰۰ تا ۵٫۰۰۰ نفر متغیر است.
    \item اعدام‌ها شامل اعضای مجاهدین، چپ‌ها (توده، فدایی، پیکار)، و حتی بی‌طرف‌ها بود.\footnote{\en{Abrahamian (1999), pp.\,209--228;} همچنین: \en{Geoffrey Robertson, The Massacre of Political Prisoners in Iran, 1988, London: Boroumand Foundation, 2010.}}
\end{itemize}

\textbf{متن فرمان خمینی (منتشرشده توسط منتظری):}

\textit{«کسانی که در زندان‌ها... بر سر موضع نفاق خود هستند، محاربند و حکمشان اعدام است... رحم بر محاربین ساده‌لوحی است.»}\footnote{متن فتوای خمینی، منتشرشده در خاطرات منتظری و تأیید‌شده توسط چندین منبع. \en{Robertson (2010), pp.\,40--55.}}
\end{warningBox}

\begin{noteBox}[منتظری و اعتراض به اعدام‌ها]
آیت‌الله حسینعلی منتظری — جانشین رسمی خمینی — یکی از معدود افرادی بود که به اعدام‌های ۱۳۶۷ اعتراض کرد. او در نامه‌هایی به خمینی نوشت:

\textit{«آقا، اعدام چند هزار نفر در چند روز... جهان ما را محکوم خواهد کرد... این با اسلام سازگار نیست.»}\footnote{نامه‌های منتظری به خمینی، منتشرشده در خاطرات منتظری، ص ۴۲۰–۴۴۵.}

نتیجهٔ اعتراض منتظری: عزل از جانشینی (فروردین ۱۳۶۸) و حصر خانگی تا پایان عمر. خمینی در نامه‌ای به منتظری نوشت: \textit{«شما ساده‌لوح هستید... منافقین شما را فریب داده‌اند.»}\footnote{صحیفهٔ امام، ج۲۱، ص ۳۳۰–۳۳۵.}
\end{noteBox}

% ────────────────────────────────────────
\section{فرایند تصفیه: گام‌به‌گام تا تمرکز کامل}
% ────────────────────────────────────────

\begin{figure}[H]
\centering
\begin{tikzpicture}[
    x=2.5cm,
    phase/.style={
        rectangle, rounded corners=4pt,
        draw=#1!70, fill=#1!10,
        text width=4cm, align=center,
        font=\small, inner sep=5pt,
        minimum height=2.5cm
    },
    arrow/.style={
        -{Stealth[length=5pt]},
        line width=1.5pt,
        #1!60
    },
    label/.style={
        font=\scriptsize\bfseries,
        color=cGray
    }
]

\node[phase=cGreen] (p1) at (0,0) {
    \textbf{مرحله ۱}\\[2pt]
    {\footnotesize بهمن ۵۷ — آبان ۵۸}\\[4pt]
    حذف لیبرال‌ها\\
    (استعفای بازرگان)\\[2pt]
    \textcolor{cGray}{\scriptsize بهانه: گروگان‌گیری}
};

\node[phase=cGold] (p2) at (2,0) {
    \textbf{مرحله ۲}\\[2pt]
    {\footnotesize آبان ۵۸ — خرداد ۶۰}\\[4pt]
    حذف ملی‌گرایان\\
    (عزل بنی‌صدر)\\[2pt]
    \textcolor{cGray}{\scriptsize بهانه: «عدم کفایت»}
};

\node[phase=cAccent] (p3) at (4,0) {
    \textbf{مرحله ۳}\\[2pt]
    {\footnotesize تیر ۶۰ — ۶۲}\\[4pt]
    سرکوب مجاهدین و چپ\\
    (اعدام‌های گسترده)\\[2pt]
    \textcolor{cGray}{\scriptsize بهانه: ترور و «نفاق»}
};

\node[phase=cConsolid] (p4) at (6,0) {
    \textbf{مرحله ۴}\\[2pt]
    {\footnotesize ۶۲ — ۶۸}\\[4pt]
    حذف منتقدان درون‌نظامی\\
    (عزل منتظری)\\[2pt]
    \textcolor{cGray}{\scriptsize بهانه: «ساده‌لوحی»}
};

\draw[arrow=cGreen] (p1) -- (p2);
\draw[arrow=cGold] (p2) -- (p3);
\draw[arrow=cAccent] (p3) -- (p4);

% برچسب بالا
\node[label, above=0.3cm of p1] {ائتلاف فراگیر};
\node[label, above=0.3cm of p2] {لیبرال‌ها رفتند};
\node[label, above=0.3cm of p3] {اپوزیسیون نابود};
\node[label, above=0.3cm of p4] {یکدست‌سازی کامل};

\end{tikzpicture}
\caption{چهار مرحلهٔ تصفیه و تمرکز قدرت (۱۳۵۷–۱۳۶۸)}
\label{fig:purge-phases}
\end{figure}

% ────────────────────────────────────────
\section{تفسیرهای رقیب: تصفیه‌ها از منظر چهار تز}
% ────────────────────────────────────────

\begin{principleBox}[چهار خوانش از فرایند تصفیه]
\begin{description}[style=nextline, font=\bfseries\color{cPrimary}]
    \item[تز ۱ (صداقت + اجبار)]
    هر مرحلهٔ تصفیه واکنش به یک تهدید واقعی بود: گروگان‌گیری → بازرگان مجبور به استعفا شد؛ ناکارآمدی بنی‌صدر → عزل ضروری بود؛ ترورهای مجاهدین → سرکوب اجتناب‌ناپذیر بود؛ اشتباهات منتظری → عزل لازم بود. خمینی هر بار «ناچار» بود.

    \textbf{نقد:} اگر چهار بار پشت سر هم «ناچار» به حذف رقیب بشوید و هر بار قدرت‌تان بیشتر شود، آیا واقعاً «ناچار» بودید یا \keyword{الگویی} وجود دارد؟

    \item[تز ۲ (خدعه)]
    تصفیه‌ها بخشی از \keyword{طرح بلندمدت} بودند: ابتدا از لیبرال‌ها استفاده کن (بازرگان)، سپس حذف‌شان کن؛ سپس از ملی‌گرایان استفاده کن (بنی‌صدر)، سپس حذف‌شان کن؛ سپس از خود مجاهدین برای ایجاد «بهانهٔ سرکوب» بهره ببر. هر مرحله بهانه‌اش مختلف بود اما نتیجه یکسان: \keyword{تمرکز بیشتر}.

    \textbf{نقد:} فرض «طرح بلندمدت» نیازمند فرض یک «ذهن تمام‌بین» است — آیا واقعاً خمینی (یا حلقهٔ درونی) تا این حد طراحی کرده بودند؟

    \item[تز ۳ (تکامل)]
    هر بحران، خمینی را یک قدم بیشتر متقاعد کرد که \keyword{تمرکز قدرت ضروری} است. ناکارآمدی بازرگان → «لیبرال‌ها نمی‌توانند»؛ شکست بنی‌صدر → «ملی‌گرایان نمی‌فهمند»؛ ترورهای مجاهدین → «رحم ضعف است»؛ اعتراض منتظری → «حتی نزدیک‌ترین‌ها خطا می‌کنند.» تکامل فکری مستمر.

    \textbf{نقد:} هر بار «تکامل» دقیقاً در جهت تمرکز بیشتر قدرت بود — آیا این «تکامل» است یا «منفعت‌طلبی»؟

    \item[تز ۴ (فشار نخبگان)]
    تصفیه‌ها محصول \keyword{رقابت جناحی} بودند: بهشتی علیه بازرگان؛ حزب ج.ا. علیه بنی‌صدر؛ سپاه علیه مجاهدین؛ خامنه‌ای و رفسنجانی علیه منتظری. خمینی در هر مرحله «داور» بود — و داوری‌اش همواره به سود جناح \keyword{وفادارتر به خودش} بود.

    \textbf{نقد:} «داوری همیشه یک‌طرفه» خودش نشانهٔ موضع‌گیری آگاهانه است — نه بی‌طرفی.
\end{description}
\end{principleBox}

% ────────────────────────────────────────
\section{مجاهدین خلق: ارزیابی انتقادی}
% ────────────────────────────────────────

\begin{noteBox}[تحفظ اخلاقی و تحلیلی]
ارزیابی نقش مجاهدین خلق نیازمند \keyword{توازن} است:
\begin{itemize}[label=\textcolor{cGold}{$\bullet$}, rightmargin=2em]
    \item عملیات‌های مسلحانهٔ مجاهدین (به‌ویژه ترورهای ۷ تیر و ۸ شهریور) قربانیان بسیاری داشت و از نظر حقوقی و اخلاقی محکوم‌کردنی است.
    \item اما واکنش نظام (اعدام‌های گسترده، شکنجه، کشتار ۱۳۶۷) نیز \keyword{نامتناسب} و ناقض حقوق بشر بود.
    \item هر دو طرف مسئولیت دارند — اما مسئولیت \keyword{حکومت} به‌دلیل برخورداری از قدرت دولتی و مسئولیت حفاظت از شهروندان، \keyword{سنگین‌تر} است.
    \item مجاهدین بعدها به عراق رفتند و در کنار صدام علیه ایران جنگیدند (عملیات فروغ جاویدان / مرصاد، ۱۳۶۷) — عملی که اعتبارشان را حتی نزد بسیاری از مخالفان نظام از بین برد.\footnote{\en{Abrahamian, Radical Islam (1989), pp.\,245--260.}}
\end{itemize}
\end{noteBox}

\subsection{آیا خشونت مجاهدین «اجتناب‌ناپذیر» بود؟}

\begin{warningBox}[پرسش ضدواقعی (\en{counterfactual})]
اگر مجاهدین در خرداد ۱۳۶۰ مبارزهٔ مسلحانه را شروع \keyword{نمی‌کردند}، آیا مسیر تمرکز قدرت متفاوت می‌بود؟

\begin{description}[style=nextline, font=\bfseries\color{cPrimary}]
    \item[احتمال اول:] بله — بدون خشونت مجاهدین، بهانهٔ سرکوب وجود نداشت و فضای سیاسی بازتر می‌ماند.

    \item[احتمال دوم:] خیر — بنی‌صدر در هر صورت عزل می‌شد (فرایند عزل قبل از ۳۰ خرداد شروع شده بود). بهانهٔ دیگری پیدا می‌شد. تمرکز قدرت «منطق ساختاری» داشت نه «بهانه‌ای».

    \item[ارزیابی:] شواهد بیشتر با احتمال دوم سازگار است — عزل بنی‌صدر \keyword{قبل از} عملیات مسلحانهٔ مجاهدین آغاز شده بود. اما خشونت مجاهدین مسلماً فرایند سرکوب را \keyword{تشدید و گسترش} داد.
\end{description}
\end{warningBox}

% ────────────────────────────────────────
\section{جدول تطبیقی نهایی: وعده‌ها در برابر واقعیت ۱۳۶۰}
% ────────────────────────────────────────

\begin{longtable}{%
    >{\centering\arraybackslash}p{0.04\textwidth}
    >{\raggedleft\arraybackslash}p{0.16\textwidth}
    >{\raggedleft\arraybackslash}p{0.35\textwidth}
    >{\raggedleft\arraybackslash}p{0.35\textwidth}}
\toprule
\rowcolor{cAccent!10}
\textbf{\#} & \textbf{محور} & \textbf{وعدهٔ پاریس (۱۳۵۷)} & \textbf{واقعیت (تا پایان ۱۳۶۰)} \\
\midrule
\endhead

۱ &
نقش خمینی &
«مقامی نخواهم داشت» &
ولی فقیه، فرمانده‌کل، عزل‌کنندهٔ رئیس‌جمهور \\
\midrule

۲ &
حکومت روحانیون &
«روحانیون حکومت نخواهند کرد» &
رئیس‌جمهور، نخست‌وزیر، رئیس قوهٔ قضاییه، رئیس مجلس — همه روحانی \\
\midrule

۳ &
آزادی احزاب &
«آزادی همه، حتی کمونیست‌ها» &
تمام احزاب غیرحکومتی ممنوع یا سرکوب‌شده \\
\midrule

۴ &
آزادی مطبوعات &
«مطبوعات آزاد» &
ده‌ها نشریه توقیف؛ فقط رسانه‌های حکومتی \\
\midrule

۵ &
حقوق متهمان &
«عدالت برای همه» &
اعدام‌های بدون وکیل، بدون استیناف \\
\midrule

۶ &
دموکراسی &
«مردم تصمیم می‌گیرند» &
ولی فقیه رئیس‌جمهور منتخب ۷۶\% را عزل کرد \\
\midrule

۷ &
رابطه با غرب &
«روابط عادی» &
قطع کامل با آمریکا، انزوای بین‌المللی \\
\midrule

۸ &
بازگشت به قم &
«به قم بر می‌گردم» &
مقر دائمی در جماران (شمال تهران) \\
\midrule

۹ &
حقوق اقلیت‌ها &
«حقوق همه محفوظ» &
سرکوب بهائیان، محدودیت سنی‌ها \\
\midrule

۱۰ &
حکومت متخصصان &
«حکومت دست متخصصان» &
تصمیم‌گیری توسط فقها و نظامیان \\

\bottomrule
\caption{جدول تطبیقی نهایی: ده محور وعده در برابر واقعیت (تا پایان ۱۳۶۰)}
\label{tab:final-comparison-1360}
\end{longtable}

\begin{warningBox}[نتیجه‌گیری اولیه]
در \keyword{تمام} ده محور، واقعیت تا پایان ۱۳۶۰ (سه سال پس از وعده‌ها) با وعده‌های پاریس در تضاد کامل قرار داشت. این تضاد صرفاً جزئی یا تاکتیکی نبود — بلکه \keyword{ساختاری و بنیادین} بود: نوع حکومتی که شکل گرفت، دقیقاً نقطهٔ مقابل نوع حکومتی بود که وعده داده شده بود.
\end{warningBox}

% ────────────────────────────────────────
\section{عزل منتظری و یکدست‌سازی نهایی (فروردین ۱۳۶۸)}
% ────────────────────────────────────────

آخرین مرحلهٔ تصفیه، عزل آیت‌الله منتظری از جانشینی خمینی بود. منتظری:

\begin{itemize}[label=\textcolor{cPrimary}{$\blacktriangleright$}, rightmargin=2em]
    \item از ۱۳۶۴ رسماً «قائم‌مقام رهبری» (جانشین) بود.
    \item به‌تدریج منتقد شد: از اعدام‌ها، از سرکوب مطبوعات، و از تمرکز قدرت.
    \item اعتراضش به کشتار ۱۳۶۷ آخرین قطره بود.
    \item خمینی در فروردین ۱۳۶۸ او را عزل کرد.\footnote{صحیفهٔ امام، ج۲۱، ص ۳۳۰–۳۴۰.}
\end{itemize}

\begin{noteBox}[اهمیت عزل منتظری]
عزل منتظری به دو دلیل مهم است:
\begin{enumerate}[label=\textcolor{cGold}{\bfseries\arabic*.}, rightmargin=2em]
    \item نشان داد حتی \keyword{نزدیک‌ترین یار} خمینی اگر منتقد شود حذف می‌شود.
    \item باعث شد پس از مرگ خمینی (خرداد ۱۳۶۸)، علی خامنه‌ای — که نه مرجع تقلید بود و نه سابقهٔ فقهی کافی داشت — ولی فقیه شود. یعنی عزل منتظری \keyword{مسیر تاریخ ایران را تغییر داد}.\footnote{\en{Moin (1999), pp.\,280--295.}}
\end{enumerate}
\end{noteBox}

% ────────────────────────────────────────
\section{جمع‌بندی فصل دوازدهم و بخش دوم}
% ────────────────────────────────────────

\begin{principleBox}[خلاصهٔ فصل ۱۲ و بخش دوم]
\begin{itemize}[label=\textcolor{cPrimary}{$\blacksquare$}, rightmargin=2em]
    \item سال ۱۳۶۰ سال تصفیه‌های نهایی بود: عزل بنی‌صدر، سرکوب مجاهدین، موج اعدام‌ها.
    \item خمینی مستقیماً در عزل بنی‌صدر و صدور فرمان اعدام‌ها نقش داشت — نه صرفاً «تأییدکننده».
    \item عملیات‌های مسلحانهٔ مجاهدین بهانهٔ مهمی برای سرکوب فراهم کرد — اما فرایند تمرکز قبل از آن شروع شده بود.
    \item کشتار ۱۳۶۷ با فرمان مستقیم خمینی انجام شد — و اعتراض منتظری به عزل او منجر شد.
    \item جدول تطبیقی نهایی نشان می‌دهد: در \keyword{تمام} ده محور، واقعیت ۱۳۶۰ نقیض وعده‌های ۱۳۵۷ است.
    \item \keyword{نتیجه‌گیری بخش دوم:} بررسی تفصیلی هفت بزنگاه تاریخی (پاریس، بختیار، قانون اساسی، فرقان، گروگان‌گیری، جنگ، تصفیه‌ها) نشان می‌دهد الگوی ثابتی وجود دارد: در هر بزنگاه، خمینی گزینه‌ای را انتخاب کرد که به تمرکز بیشتر قدرت منجر شد. این الگو با تزهای دوم (خدعه) و چهارم (فشار نخبگان) بیشترین سازگاری را دارد — هرچند تز سوم (تکامل) نیز بخشی از تبیین را فراهم می‌کند. تز اول (صداقت + اجبار) ضعیف‌ترین عملکرد را در مجموع دارد.
\end{itemize}
\end{principleBox}

\vspace{1cm}

\begin{center}
\begin{tikzpicture}
    \draw[cAccent, line width=0.5pt] (0,0) -- (10,0);
    \node[circle, fill=cAccent, inner sep=2pt] at (5,0) {};
\end{tikzpicture}

\vspace{0.5cm}
{\small\color{cGray} — پایان \textbf{بخش دوم} (سیر حوادث و تحلیل تطبیقی) —}

\vspace{0.3cm}
{\small\color{cGray} ادامهٔ کتاب در \textbf{بخش سوم} (\en{Segment 4}): فصل‌های ۱۳ تا ۱۶}

\vspace{0.2cm}
{\small\color{cGray} تحلیل تطبیقی نهایی، خوانش‌های رقیب در بوتهٔ نقد، وزن‌دهی شواهد، و نتیجه‌گیری}
\end{center}